\chapter*{前言}
\addcontentsline{toc}{chapter}{前言}

\begin{notebox}
\textbf{写给翻开这本书的你}

如果你是一名水利工程专业的大三学生，你可能正在思考一个问题：学水利的，为什么要学写代码？

答案就藏在你身边。打开手机查看天气预报中的降水预警，浏览新闻里水库实时水位的三维动态图，或者在实习中看到监测站大屏上跳动的数据曲线——这些背后，都是软件平台在支撑。今天的水利工程师，不仅要懂水力学和结构力学，还需要理解这些平台是如何设计和构建的。这正是本书希望带给你的能力。
\end{notebox}

\section*{为什么写这本书}

近年来，"数字孪生水利""智慧水利"已从政策文件走进了工程实践。《国家水网建设规划纲要》明确提出建设数字孪生水利体系的目标，各流域管理机构和水利工程单位对兼具水利专业背景和软件开发能力的复合型人才需求急剧增长。

然而在教学中，我们面临一个尴尬的现实：计算机专业的软件工程教材讲得深，但案例全是电商、社交，与水利毫无关系；水利信息化教材介绍了不少概念，但很少教学生怎样动手搭建一个真正能跑起来的平台。本教材的初衷，就是填补这个空白——用水利工程的真实场景来串联软件开发的完整流程，让你学到的每一行代码都与专业相关。

\section*{这本书讲什么}

全书九章，按照\textbf{“打基础 $\rightarrow$ 学技术 $\rightarrow$ 做项目”}的思路组织，分为三个部分。

\begin{importantbox}
\textbf{第一部分\quad 基础篇（第1--3章）}

从智慧水利的概念出发，建立对平台整体架构的认知。然后进入软件工程的核心内容：需求怎么分析、系统怎么设计、架构怎么选型。这部分是后续所有开发工作的理论根基。
\end{importantbox}

\begin{tipbox}
\textbf{第二部分\quad 技术篇（第4--7章）}

四章对应平台开发与可视化表达的核心技术链：
\begin{itemize}
\item \textbf{前端}（第4章）：HTML/CSS/JavaScript基础，Vue.js框架开发，重点是如何构建水情数据展示界面
\item \textbf{后端}（第5章）：Spring Boot和Flask框架，RESTful API设计，数据持久化与安全认证
\item \textbf{三维场景}（第6章）：WebGL/Three.js三维渲染，GIS地图服务，BIM建模与数字孪生技术
\item \textbf{观测数据展示}（第7章）：监测数据面向三维场景的组织、可视化表达、交互与联动
\end{itemize}
每个技术点都围绕水利监测平台的实际需求展开，代码示例均取自水利业务场景。
\end{tipbox}

\begin{notebox}
\textbf{第三部分\quad 应用篇（第8--9章）}

前两部分学到的方法与技术，在这里汇聚成完整的项目实践。第8章以\textbf{水利工程安全监测平台}为起点，从业务需求、总体架构和数据模型出发，完成前后端、三维场景、智能分析、部署与验收的完整链路；并在8.5节进一步演进到数字孪生水利平台，以流域、水网和工程三类入口贯通“预报、预警、预演、预案”。第9章回顾全书，并给出技术展望与后续学习指引。
\end{notebox}

\section*{怎么使用这本书}

\textbf{先说最重要的一条：一定要动手写代码。}

软件开发不是看会的，是写会的。每章末尾都有思考题和练习题，思考题帮你梳理概念，练习题要求你实际编码。请认真对待它们。

以下是几条具体建议：

\begin{enumerate}
\item \textbf{按顺序学习。}各章内容前后衔接，第4--7章的开发与可视化技术依赖第1--3章的方法基础，第8章的综合实践又建立在第4--7章的技术基础上。不建议跳章阅读。

\item \textbf{搭建本地开发环境。}学习前端章节时，准备好Node.js和VS Code；学习后端章节时，安装好Python或Java开发环境。在自己的电脑上运行书中的示例，比看十遍书都有效。

\item \textbf{组队做课程设计。}如果条件允许，建议3--4人组成开发小组，分工协作完成一个小型水利监测平台。前端、后端、三维场景各有分工，在协作中体会真实项目的开发流程。

\item \textbf{善用搜索引擎和技术社区。}书中不可能涵盖所有细节。遇到问题时，学会查阅官方文档、在Stack Overflow或CSDN上搜索解决方案，这本身就是软件工程师的核心技能之一。
\end{enumerate}

\section*{推荐教学组织方式}

为了便于任课教师直接采用，本书建议按以下三种方式组织教学：

\begin{itemize}
\item \textbf{32学时方案：}以第1--6章为主，突出智慧水利平台的总体认知、需求分析、架构设计以及前后端和三维场景的核心概念，适合作为专业方向课或选修课。
\item \textbf{48学时方案：}完整讲授第1--8章，并结合章末交付物组织阶段作业，使学生最终完成一个最小可运行的水利监测平台原型。
\item \textbf{课程设计方案：}以第8章贯穿案例为中心，要求学生分组完成需求说明、架构设计、前端页面、后端接口和三维场景集成，可与毕业设计或综合实训衔接。
\end{itemize}

\section*{适用对象}

\textbf{本书以水利工程及相关专业高年级本科生为主要读者}，兼顾研究生，以及有志于从事水利信息化、智慧水利平台开发的技术人员。

\textbf{先修知识依赖}：学习本书前，你需要具备以下基础：

\begin{itemize}
\item 至少学过一门编程语言（C、Python或Java均可）
\item 了解基本的数据结构概念（数组、链表、树等）
\item 对水利工程有初步认识（水力学、水文学等先修课程）
\item 了解基本的数据库概念与SQL语言（推荐但不强制）
\end{itemize}

本书也可作为水利行业技术人员了解平台开发技术的参考读物。

\section*{致谢}

本教材的编写得益于多年的教学实践积累和工程项目经验，参考了大量国内外文献和优秀的开源项目。感谢在教材编写过程中提出宝贵意见的各位同事和同学，感谢审稿专家的悉心指导。

智慧水利技术仍在快速发展之中，书中难免存在不足和疏漏。欢迎读者通过各种渠道提出批评和建议，帮助我们在后续版本中不断完善。

\vspace{2em}
\begin{flushright}
编者\\
2026年8月
\end{flushright}
